\section{The differential of the link period}\label{sec:hodge}

We first compute the coefficient of $\varepsilon$ in a period over the
dual numbers. Its component of Hodge degree two is contraction by the
Kodaira--Spencer class. Logarithmic forms on a resolution identify
this component in the Hodge
filtration of the link. To compare it with the period itself we construct
a holomorphic primitive for every holomorphic three-form on the punctured
toric germ.

\subsection{Logarithmic forms at a toric point}
Let $(S,p)$ be an isolated Gorenstein toric threefold singularity. Take
a sufficiently small contractible Stein representative and write
$S^\circ=S\setminus\{p\}$. On a good resolution $M\to S$, let $D$ be
the reduced exceptional divisor. We take representatives for which $M$
retracts onto $D$ and $M\setminus D=S^\circ$ has the homotopy type of
the link $L$. Put $\mathcal F^a=\Omega_M^a(\log D)$. Unless indicated
otherwise, topological cohomology has complex coefficients.

Steenbrink's toric vanishing theorem gives
\begin{equation}\label{hg:vanishing}
 H^q(M,\mathcal F^a(-D))=0\qquad (q>0)
\end{equation}
on any good resolution \cite[\S2, Theorem~2]{Ste}. Local duality
identifies the dimension of this group with that of
$H_D^{3-q}(M,\mathcal F^{3-a})$ \cite[Definition~1.4]{FL}.
In particular $H^1_D(M,\mathcal F^2)=H^2_D(M,\mathcal F^2)=0$.
The exact sequence for cohomology with support and the restriction
sequence for $D$ now give isomorphisms
\begin{equation}\label{hg:restrictions}
 H^1(S^\circ,\Omega^2_{S^\circ})
 \xleftarrow{\ \sim\ } H^1(M,\mathcal F^2)
 \xrightarrow{\ \sim\ } H^1(D,\mathcal F^2|_D)
 =\Gr_F^2 H^3(L).
\end{equation}
The last equality comes from the logarithmic complex
$\mathcal F^\bullet|_D$ for the link, whose Hodge spectral sequence
degenerates at $E_1$ \cite[\S3]{Ste}.

In fact $H^3(L)$ is pure of type $(2,2)$. Toric singularities are
rational, so $\Gr_F^0H^2(L)=\Gr_F^2H^2(L)=0$ by
\cite[\S3, Lemma~2]{Ste}. The logarithmic complex has no other
possible Hodge degree in cohomological degree two except one.
Hodge symmetry on each weight quotient therefore leaves only type
$(1,1)$ in $H^2(L)$. Link duality identifies $H^3(L,\Q)$ with
$H^2(L,\Q)^\vee(-3)$ as a mixed Hodge structure \cite[\S3]{Ste}.
In particular $\Gr_F^2H^3(L)=H^3(L)$ canonically.

Choose a generator $\omega$ of the dualizing sheaf on $S$. The
depth-three comparison identifies the first-order deformation space
$T^1_{S,p}$ with $H^1(S^\circ,\Theta_{S^\circ})$
\cite[Lemma~1.17 and Remark~4.2]{FL}. Here the usual cotangent-complex
tangent group agrees with $\Ext^1(\Omega^1_S,\OO_S)$ because the negative
cohomology sheaves of the cotangent complex have finite support and
grade three in $\OO_S$. Contraction with $\omega$ on
$S^\circ$, followed by \eqref{hg:restrictions}, defines an isomorphism
\begin{equation}\label{hg:lambda}
 \lambda_\omega:T^1_{S,p}\longrightarrow H^3(L).
\end{equation}

\begin{proposition}\label{prop:variation}
Let $A=\C[\varepsilon]/(\varepsilon^2)$. On a deformation of $(S,p)$
over $A$ with Kodaira--Spencer class $v$, every relative volume form
lifting $\omega$ has marked link class
$\varepsilon\lambda_\omega(v)$. In particular the differential of the
period of the homogeneous toric generator is an isomorphism from
$W_Q$ to $H^3(L)$.
\end{proposition}

\begin{proof}
We shall use the exactness of holomorphic three-forms on $S^\circ$
after shrinking $S$. Write such a form as $f\omega_0$, where $\omega_0$
is the toric generator and $f$ extends across $p$ by normality.
An integral cocharacter in the interior of the cone contracts $S$ to
$p$. Its Euler field $e$ has positive weights on the nonconstant
semigroup generators, and $\omega_0$ has positive weight $a$, the
pairing of the cocharacter with the Gorenstein degree. On a small
invariant polydisc for the contracting real scalars $0<s\leq1$, set
\[
 F(x)=\int_0^1 s^{a-1} f(s\cdot x)\,ds .
\]
In homogeneous ambient coordinates the integral divides the coefficient
of each monomial of weight $m$ by $a+m$, so $F$ is holomorphic. Thus
$(e+a)F=f$ and Cartan's formula gives
$d\iota_e(F\omega_0)=f\omega_0$. In particular choose
$\Gamma$ on $S^\circ$ with $d\Gamma=\omega$.

Trivialize the first-order deformation on a Stein cover of $S^\circ$.
Use the convention that transport of forms from chart $j$ to chart $i$
is $1-\varepsilon\mathcal L_{\kappa_{ij}}$ for the
Kodaira--Spencer cocycle $\kappa$. Write the lifted form in chart $i$
as $\omega+\varepsilon\eta_i$ and put
$\beta_{ij}=\iota_{\kappa_{ij}}\omega$. The gluing identities give
$\delta\eta=d\beta$ and $\delta\beta=0$. With total differential
$\delta+(-1)^q d$ in \v Cech degree $q$, the pair $(\eta,\beta)$
represents the coefficient of $\varepsilon$ in the marked class.
Indeed lift $\Gamma$ by its central expression on each chart.
Subtracting the total coboundary of this lift from the relative volume
form leaves $\varepsilon(\eta,\mathcal L_\kappa\Gamma)$ on the central
complex. The second entry differs from $\beta$ by
$d\iota_\kappa\Gamma$, a total coboundary since
$\delta\iota_\kappa\Gamma=0$. The relative de Rham resolution of the
constant sheaf $A$ identifies multiplication by $\varepsilon$ in this
calculation with the coefficient inclusion $\C\to A$,
$c\mapsto\varepsilon c$.

Lift $[\beta]$ by \eqref{hg:restrictions} and choose a logarithmic
one-cocycle $b$ representing the lift. We claim that $db$ is a
\v Cech coboundary on $M$. By \eqref{hg:vanishing} for $a=3$ the
restriction
$H^1(M,\mathcal F^3)\to H^1(D,\mathcal F^3|_D)$ is an
isomorphism. There the class of $db$ is the first differential of
the link Hodge spectral sequence on $[b|_D]$, hence zero. We may
therefore choose a zero-cochain $c$ of logarithmic three-forms with
$\delta c=db$. The total cocycle $(c,b)$ on $M$ restricts to a
class on the link whose Hodge degree-two component is $[b|_D]$.
As $F^2H^3(L)=H^3(L)$ and $F^3H^3(L)=0$, that component determines
the class and gives $\lambda_\omega(v)$.

The logarithmic comparison identifies restriction of $(c,b)$ to
$S^\circ$ and to the link complex under $H^3(S^\circ)\simeq H^3(L)$
\cite[\S1.1]{FL}. On $S^\circ$ choose a zero-cochain $u$ of
two-forms with $\delta u=\beta-b$, refining the covers as necessary.
The difference of $(\eta,\beta)$ and $(c,b)$, modulo the total
coboundary of $u$, is the global holomorphic three-form
$\eta-c-du$. Its primitive constructed above gives zero in $H^3(L)$.
This proves the formula for every lift of $\omega$. For $S$ the cone
of $Q$, the Kodaira--Spencer isomorphism $W_Q\simeq T^1_{S,p}$
of \cite[\S6.2 and the corollary in \S7.2]{Alt} gives the final assertion.
\end{proof}

\subsection{Restriction from the smooth locus}
Let $X$ be a normal projective threefold with isolated Gorenstein toric
singularities. Choose a global volume form $\omega$ trivializing its
dualizing sheaf. Write
$U=X_{\mathrm{reg}}$. Let $Y\to X$ be a projective good resolution
which is an isomorphism over $U$, with reduced exceptional divisor
$E=\coprod_p E_p$. Orient each $L_p$ as the boundary of a small
resolution neighborhood of $E_p$.

\begin{proposition}\label{prop:tangent}
Composing tangent localization with the local period differentials for
$\omega$ and link Poincar\'e duality gives an exact sequence
\[
 0\longrightarrow H^1(X,\Theta_X)
 \longrightarrow T_0\Def(X)
 \longrightarrow K\longrightarrow 0 .
\]
Under the same duality the image of
$H^3(U)\to\bigoplus_p H^3(L_p)$ is $K$.
\end{proposition}

\begin{proof}
Put $\mathcal F=\Omega_Y^2(\log E)$ and let $Y_p$ be the inverse
image of a small Stein representative at $p$. Excision identifies
$H^i_E(Y,\mathcal F)$ with the direct sum of the corresponding
groups on $Y_p$. Their vanishing for $i=1,2$ proves that
$H^1(Y,\mathcal F)\to H^1(U,\Omega_U^2)$ is an isomorphism.
The natural depth-three comparison
$T_0\Def(X)\simeq H^1(U,\Theta_U)$ intertwines tangent localization
with restriction to the punctured representatives
\cite[Lemma~4.1 and Remark~4.2]{FL}. Under contraction with
$\omega$ and the inverse restriction isomorphisms, localization
therefore becomes the map
\[
 H^1(Y,\mathcal F)\longrightarrow\bigoplus_p H^1(Y_p,\mathcal F|_{Y_p}).
\]

For $j:U\hookrightarrow Y$, restriction of $Rj_*\Q_U$ to $E$ gives the
map $r:H^3(U)\to\bigoplus_p H^3(L_p)$. Its de Rham realization is
restriction of $\Omega_Y^\bullet(\log E)$ to $E$, with the filtration
by form degree \cite[\S1.1]{FL}. In particular $r$ is a morphism of
mixed Hodge structures. By logarithmic Hodge degeneration for the
smooth projective pair $(Y,E)$ \cite[Corollary~3.2.13(ii)]{Del}, the
source of the last display is $\Gr_F^2H^3(U)$. After the local
identifications \eqref{hg:restrictions}, the displayed map is
$\Gr_F^2r$. Its target has only Hodge type $(2,2)$, so strictness
\cite[Theorem~2.3.5(iii)]{Del} gives
$\im(\Gr_F^2r)=\im r$. Proposition~\ref{prop:variation} identifies
this image with that of the differential of the marked periods.

The Mayer--Vietoris sequence for $Y=U\cup\bigcup_pY_p$ contains
\[
 H^3(U)\oplus\bigoplus_p H^3(E_p)
 \longrightarrow \bigoplus_p H^3(L_p)
 \xrightarrow{\partial} H^4(Y).
\]
The arrow $H^3(E_p)\to H^3(L_p)$ respects mixed Hodge structures
\cite[(1.4)]{FL}. Its source has weights at most three since $E_p$ is
proper, whereas the target is pure of weight four. This arrow
vanishes, giving $\im r=\ker\partial$. The boundary map is the sum
of the Gysin maps of the oriented links, as follows by excision and the
Thom isomorphism on their collars. Poincar\'e duality thus identifies
its kernel with the kernel of $\bigoplus_p H_2(L_p)\to H_2(Y)$.
This proves both image assertions. Finally, the local-to-global
spectral sequence for the cotangent complex identifies the kernel of
tangent localization with $H^1(X,\Theta_X)$. Each local period
differential is an isomorphism by Proposition~\ref{prop:variation},
so this is also the kernel in the asserted sequence.
\end{proof}
